% Appendix C -- the detailed adjudication of related work, moved out of
% section 7 (2026-09-26) for the page limit. The body keeps every claim this
% paper's positioning rests on; what is here is the evidence for those claims,
% which a reviewer checking our reading of another paper will want and a reader
% following the argument will not.

\section{Related work in detail}
\label{app:related}

\paragraph{Where Pie's advantage comes from, and where it stops.} Pie's
headline is $1.9\times$ over vLLM, and the regime that produces it is stated
plainly in their \S5.3: they hold the physical KV cache fixed and sweep it from
3 to 12\,GB, and report that Pie ``outperforms vLLM at smaller GPU KV cache
sizes and achieve similar performance at the largest memory capacity
(12\,GB)'', because ``when GPU memory capacity is small, vLLM suffers from
significant recomputation''. As memory grows, ``the improvement ratio gradually
decreases until Pie and vLLM perform similarly''.

The gain is therefore \emph{avoided recomputation} in a memory-starved engine,
not additional capacity converted into work. That distinction matters for what
each system can claim. A device large enough that the baseline stops preempting
is a device on which Pie's reported advantage has already converged to zero ---
and their entire sweep, 3 to 12\,GB of KV, sits below the operating point we
study, where a 32B model on a 96\,GB device leaves $\Ckv \approx 22$--28\,GB.

\paragraph{A consequence of the layer unit: the return leg.} Choosing the layer
as the unit of management has a second effect, on traffic rather than on
policy. In Pie a layer's KV \emph{lives} on one side or the other, and the
GPU-side buffer it occupies is needed by the next layer arriving, so a layer
that has just been computed --- and has therefore just had this token's $K,V$
appended to it --- must be written back in full. Their swap controller says so
directly: it initiates ``a swap-out for the `coldest layer', defined as the last
layer that finished computation''. Every offloaded byte therefore crosses the
interconnect \textbf{twice per generated token}, and the return leg is the
slower direction (371 vs 419\,GB/s by their own measurement).

Our rolling region is not the home of the data but a cache of it: the host
mirror stays authoritative and each step's new $K,V$ is written to it directly,
so a slot can be overwritten without persisting anything. DirectKV's kernel has
the same property, appending only the newly projected $K,V$ to CPU memory. The
resulting difference in per-token traffic is not marginal:

\begin{table*}[t]
\centering
\small
\begin{tabular}{lll}
\toprule
per generated token & to device & to host \\
\midrule
Pie~\cite{pie} & offloaded KV & \textbf{the same KV again} \\
DirectKV~\cite{directkv} & all KV & one token's $K,V$ \\
\textbf{this paper} & selected requests' KV & one token's $K,V$ \\
\bottomrule
\end{tabular}
\caption{Per byte of offloaded KV, a design whose offloaded unit is \emph{owned}
by one side must return it; one whose device-side copy is a \emph{cache} need
not. Pie moves two bytes where the other two move one, and its return leg uses
the slower direction of the link.}
\label{tab:traffic}
\end{table*}

Table~\ref{tab:traffic} states the consequence. This is not a criticism of
Pie's design so much as a property of its unit:
a whole layer's cache cannot be left dirty on the device when the buffer is
reclaimed for the next layer. It does mean that comparisons of interconnect
pressure between these systems must count both directions.

\paragraph{On comparing directly.} Pie has no public artifact --- its paper
lists no repository, and we found none --- so a head-to-head against their
implementation is not available to us. Nor would it isolate what we want:
running two engines compares implementations, whereas the question here is
whether \emph{choosing} which requests to offload beats spreading the same
interconnect budget uniformly. \S\ref{sec:eval:ablation} therefore holds the
mechanism fixed and varies the policy, and reports Pie's published figures
alongside rather than reproducing them.


% tab:endpoints moved here from section 1 (2026-09-27): the body cannot afford
% two floats in a 1.25pp introduction, and this is the paragraph that argues
% the table's point at length.
\begin{table}[t]
\centering
\small
\begin{tabular}{lccr}
\toprule
\emph{Projected} & HBM & link & batch \\
\midrule
all-HBM (vLLM, SGLang)     & 26.5\,GB      & \textbf{0}    & 12 \\
all-host (uniform offload) & \textbf{0}    & 14.2\,GB      & 6 \\
\textbf{both (this paper)} & 26.5\,GB      & 14.2\,GB      & \textbf{18} \\
\bottomrule
\end{tabular}
\caption{KV read per step, and the batch it supports, for Qwen2.5-32B fp16 at 8K
on a GH200-96; each single-memory design leaves one KV path unused.
\textbf{Projections} from \texttt{model/killer\_v2.py}: $W = 65.5$\,GB measured
less a 4\,GB device reserve gives $\Ckv = 26.5$\,GB, $\btok = 256$\,KiB,
$\mathrm{ITL}_0 = 33.9$\,ms from Eq.~\ref{eq:step}, and $\Xlink = 419$\,GB/s
measured on this part~\cite{pie} rather than the 450 spec. These count
\emph{capacity} --- how many requests fit, not how fast they are served, a
difference \S\ref{sec:model:omits} shows is decisive. All three rows share one
$\mathrm{ITL}_0$ to stay comparable, which flatters the all-host row: computed
self-consistently its step would be weights-only and its batch 4. The interior
point is unaffected.}
\label{tab:endpoints}
\end{table}

\paragraph{The endpoint argument.} DirectKV \S6: ``During prefill, the manager
stores the initial $K,V$ entries from each layer in CPU memory. During decode,
the GPU reuses these cached entries directly, while only appending new $K,V$
for the latest token'' --- and its Alg.~3 appends that new $K,V$ to CPU memory
as well. \textbf{It places 100\% of KV on the host, for every request, always.}
That is why it needs no scheduler: placement is static and total. Stock vLLM
and SGLang are the opposite endpoint, at 0\%. Pie is the one system that is
genuinely interior, and it reaches the interior by splitting \emph{layers}
uniformly across all requests rather than by choosing requests --- so the
per-request axis is unoccupied, and the optimum along it is interior
(Table~\ref{tab:endpoints}). Note that the all-host endpoint leaves the device
memory not holding weights entirely idle, which is why it loses to an
HBM-resident baseline whenever that baseline fits --- the ordering DirectKV
reports on Llama-3.1-8B, and one our gain law is consistent with
(\S\ref{sec:model:directkv}).

A natural objection to reading DirectKV \S8.1 as leaving the policy dimension
open is that the vLLM and SGLang mechanisms it names --- continuous batching,
prefix reuse, eviction, hierarchical KV caching --- already constitute KV
policy. They do not, for our purposes: each decides \emph{reuse} and
\emph{preemption}, that is, the residency of blocks belonging to a request that
is \emph{not currently decoding}. None of them decides the tier of an
\emph{actively decoding} request. Several systems do \emph{decode} from
host-resident KV --- FlexGen, Pie, NEO and OmniServe among them --- so the state
is not new; what is unoccupied is choosing, per request, which side of the
memory such a request lives on.

\paragraph{Policy, at other granularities.} OmniServe~\cite{omniserve} is the
closest prior work on policy: it forms mixed GPU/host batches and admits
host-tier work per layer under an SLO. The regime differs on three axes that
change the problem. It targets PCIe-attached A100s with CPU-only hosts over
RoCE; attention for host-tier requests is computed \emph{on the CPU}, so the
trade is compute-for-capacity rather than bandwidth-for-capacity; and tiering
is driven by an \emph{externally supplied} latency-sensitive/best-effort
service class. Ours assumes a coherent interconnect, computes all attention on
the GPU, and infers tier from observable request properties, so it applies in a
single-class deployment where no such label exists.

Fang et al.~\cite{calkv} --- whose title also claims the phrase ``heterogeneous
memory system'', so the distinction is worth drawing sharply --- give a formal
placement model and a
simulated-annealing oracle bound on aggregate bandwidth. It is simulation only,
on synthetically shrunk HBM, at \emph{token} granularity by importance, and it
requires perfect foresight; its gains are contingent on 60--80\% attention
sparsity, and it closes by explicitly calling for runtime policies. Its own
results show migration overhead dominating below roughly 60\% sparsity. We
make the testable counter-claim that \emph{request}-granularity migration
avoids that thrash, because whole requests move once at controlled moments
rather than pages chasing shifting token importance.

\paragraph{Reducing the KV instead of relocating it.} ShadowKV~\cite{shadowkv}
keeps a low-rank $K$ on the GPU, offloads $V$, and reconstructs sparsely,
reporting support for a $6\times$ larger batch and $3.04\times$ throughput; KV
quantization, Quest and H2O likewise shrink what must be stored. These change
the model's output, and they compose with a memory system rather than competing
with it: a smaller KV scales both the device-resident cohort and the affordable
host cohort, leaving the ratio between them --- and hence the value of a second
tier --- unchanged. That argument is analytical and we have not measured the
composition; it is a prediction of Eq.~\ref{eq:cancel}, and a cheap one to
falsify.


\paragraph{Is decode serving memory-bound at all?} NanoFlow~\cite{nanoflow}
argues the opposite of \S\ref{sec:background}: that end-to-end serving becomes
\emph{compute}-bound for most common workloads once operations of different
resource classes are overlapped within a device, and it reports $1.91\times$ over
prior engines by doing so. The claim is not in conflict with ours, but the boundary
between them is exactly where this paper lives. NanoFlow raises utilisation at a
\emph{given} batch by filling idle units; our premise is that the batch itself is
capped by KV capacity before those units saturate. The two compose in principle ---
a larger batch gives intra-device overlap more to work with --- and they collide in
practice if the arithmetic recovered by overlapping pushes the step into the
compute-bound regime, because then admitting a host-tier request no longer buys
throughput. Our own measurement of the margin is in
Appendix~\ref{app:params}: at batch 142 for an 8B model at 4K --- far above the
capacity limit on any device we consider --- decode-attention's arithmetic
intensity has crossed the hardware balance point, yet the step is still
memory-bound by more than $6\times$ ($46.6$\,ms of memory time against $7.4$ of
compute, taking the memory term from Eq.~\ref{eq:step} rather than from the
single-efficiency roofline \S\ref{sec:model:modeldep} refutes). So the two premises
are not in conflict at the batch sizes a capacity limit permits. On a device where
that margin closes, NanoFlow's premise holds and ours does not.

\paragraph{Moving the request instead of its memory.} Llumnix~\cite{llumnix}
reschedules requests \emph{across} model instances by live-migrating them and their
KV, to balance load, reduce fragmentation and honour priorities. It is the natural
alternative to a second memory tier: when a request will not fit here, one can move
it to another GPU rather than move its KV to the host. The choice between them is
economic rather than technical --- migration needs a second instance with room,
which is a provisioning decision, while tiering needs only host DRAM, which is
already installed and idle. They also compose cleanly, since a tier decision is
local and a migration decision global, and a scheduler holding both has a strictly
larger action set than one holding either. We do not evaluate that combination.

\paragraph{Architectures that shrink the cache.} The gate of
\S\ref{sec:model:gate} opens when weights consume enough of the device that little
is left for KV, and architectural changes to attention move that boundary directly.
Multi-head latent attention~\cite{deepseekv2} compresses $K$ and $V$ into a
low-rank latent, cutting KV bytes per token by roughly an order of magnitude, and
sliding-window and hybrid-attention designs bound the cache by construction rather
than by context. By Eq.~\ref{eq:cancel} a uniform reduction in $\btok$ scales both
cohorts and leaves their ratio unchanged, so shrinking KV does not by itself close
the gate. What closes it is the \emph{second} gate: a smaller cache raises
$B_{\mathrm{cap}}$, so the fixed per-step cost is amortised over more requests but
the capacity gain it must beat is unchanged in ratio and smaller in absolute
requests.

\paragraph{Adjacent memory management.}
InfiniGen~\cite{infinigen} attacks the same per-step transfer from the other
end, speculating which tokens matter for the next layer and prefetching only
those; it reduces the bytes a host-tier request must move, where we decide how
many requests move bytes at all, and the two are complementary.
CachedAttention~\cite{cachedattention} keeps the KV of \emph{inactive}
conversation sessions in a hierarchical store so a later turn skips prefill ---
the case \S\ref{sec:background} distinguishes from ours, and the reason our
novelty claim is scoped to actively decoding requests rather than to
host-resident KV in general. vAttention~\cite{vattention} changes how KV is
\emph{allocated}, decoupling virtual from physical memory instead of paging in
user space; it is orthogonal to placement, and a host tier built on it would
face the same admission question we pose here.



\subsection*{Composition with a zero-copy kernel}

\textbf{The staging buffer they remove is the thing that makes hiding
possible.} Their kernel reads host memory \emph{during} the attention that
needs it, so the transfer is inside the operation rather than alongside it;
ours prefetches layer $\ell{+}1$ into a rolling device buffer while layer
$\ell$ computes. The two designs therefore trade in opposite directions. They
pay no device memory and get no overlap; we pay one $k$-window of device memory
and hide about $90\%$ of the transfer behind the previous layer's compute ---
a real window, though a fixed per-step cost dwarfs what escapes it. Neither dominates on paper, and
settling it needs both mechanisms on one machine, which requires their
SM90a-only kernel and is therefore future work rather than a result we report. It also means a host-tier
request cannot simply be moved onto their kernel and remain what we call a
rider: without a prefetch window there is nothing left to hide, and such a
request pays the full transfer inside its own step.\footnote{Their architecture
would admit a mixture --- their \texttt{forward\_decode} already dispatches two
attention implementations over disjoint subsets of one batch --- but the
resulting system is a third design point, selective admission \emph{without}
hiding, not ours built on theirs.}

What their \S8.1 does establish is that the policy dimension is untouched:
``DirectKV does not change scheduling, batching, or KV-management policies;
instead, it provides an efficient execution path for CPU-resident KV
blocks\ldots while leaving the scheduler, model execution graph, and frontend
API unchanged.'' Whichever mechanism proves faster, the question of \emph{how
many} requests to serve from the host is left open by it, and that is where
this paper sits.

It is also worth recording that DirectKV's own analysis agrees with our
measured PCIe result: they conclude that ``on PCIe, DirectKV should be viewed
primarily as a capacity-extension mechanism\ldots its performance advantage
remains bounded by CPU--GPU interconnect bandwidth,'' and report NVLink-C2C
reducing attention latency by up to $4.2\times$ relative to PCIe. Our
\S\ref{sec:eval} finds the same boundary from the scheduling side. Their
evaluation covers Llama-3.1-8B and OPT-13B/30B on a GH200-96, so it does not
overlap the models we evaluate, and it does not include a 70B-class
configuration --- which our capacity model identifies as the most favourable
point of all. We do not evaluate one either (\S\ref{sec:eval:setup} says why: at
a capacity-limited batch of about seven, one request is 14\% and no result could
separate an effect from a rounding artifact), so the most favourable point of
the model's own curve is measured by nobody yet.


\subsection*{The two axes, and moving compute instead}

\paragraph{Moving compute instead of data.} NEO~\cite{neo} offloads attention
\emph{compute} to the CPU to grow the batch --- the same problem statement as
ours with the opposite mechanism. Its own scaling is the argument against it:
gains decay from $7.5\times$ on a T4 to 14\% on an H100, because a faster GPU
makes the CPU's arithmetic relatively worse. Streaming KV to the GPU depends on
a different ratio, $\Xlink/\BW$, and it would be convenient but dishonest to
present that as simply improving with hardware. Across a generation, HBM
bandwidth and capacity have both grown faster than host-to-device bandwidth, so
$\Xlink/\BW$ falls and $\HBM/\Ckv$ falls with the extra capacity --- both
factors of Eq.~\ref{eq:gain-exact} trend \emph{against} us, which is the same
structural decay we note in NEO. What would reverse it is a link that scales
with HBM rather than with PCIe generations, which is precisely the bet a
coherent interconnect represents. \S\ref{sec:discussion} states our own
anti-scaling in the same terms.

Appendix~\ref{app:related} gives our detailed reading of Pie's own sweep, the
endpoint argument from DirectKV's text, and the placement of the adjacent
memory-management work.

\paragraph{Why the request axis?} The two axes are not orthogonal, and it would
be convenient but wrong to say so: a layer split and a request split spend the
\emph{same} interconnect budget, and in the fluid model of \S\ref{sec:model} a
layer split with the right $m/n$ fills both byte budgets and reaches the same
interior point. They are substitutes, and the case for the request axis has to
be made on the terms the fluid model omits.

Three of those favour it. A host-tier request holds \emph{zero} device blocks, so
admitting one requires no allocation at all and the device pool is untouched;
under a layer split every request retains $(n-m)/n$ of its cache on the device,
so capacity is freed only in proportion and never for a whole request. Movement
happens once per request, at a moment the scheduler chooses, rather than once per
layer per token --- which is what removes Pie's return leg. And only selected
requests pay link traffic, so a load spike can be absorbed by moving a few
requests rather than by degrading all of them.

One favours the layer axis, and we concede it: a layer split has no integer
quantization. It can relocate $1/n$ of a cache where we must relocate a whole
request, so at long context, where $B_{\mathrm{cap}}$ and $B_{\mathrm{off}}$ are
small integers (\S\ref{sec:model:omits}), the layer axis reaches operating points
the request axis cannot express. A system wanting both would use a request split
for capacity and a layer split for the remainder; we have not built that, and it
is the more interesting composition than the one usually asserted.

We adopt one of Pie's measurements --- Pie reports GH200 NVLink-C2C at
419\,GB/s CPU$\to$GPU and 371\,GB/s GPU$\to$CPU, an independent confirmation
that the widely quoted 900\,GB/s is a bidirectional figure and that a
host-to-device KV stream should be budgeted well below it.

FlexGen~\cite{flexgen} is the earlier generation of the same family, placing KV
across GPU, CPU and disk by a linear-programming search with compression; it is
offline-profiled rather than adaptive. Both pay the staging cost that a
zero-copy kernel removes, and neither decides a tier per request.

DAK~\cite{dak} pushes host-resident tensors to shared memory via the Tensor
Memory Accelerator (TMA) and chooses
per-\emph{operation} offload ratios by greedy search --- but takes a
\emph{global} offload ratio as an input, which is precisely the quantity a
policy has to produce.
